\documentclass[../../../../main.tex]{subfiles}
\begin{document}

\begin{flushright}
\footnotesize\textit{【自動文字起こし・要確認】}
\end{flushright}

[解]
(イ) まず、2正数 $a, b$ に対し、$a+b \geqq 2\sqrt{ab}$ であることを示す。両辺正から2乗して良く、
\[
a+b \geqq 2\sqrt{ab} \iff (a-b)^2 \geqq 0
\]
より、示された。$a = a_1+a_2$, $b = a_3+a_4$ として ($\because a_k > 0$)
\[
\frac{1}{2}(a_1+a_2+a_3+a_4) \geqq \sqrt{(a_1+a_2)(a_3+a_4)} \quad \cdots \text{①}
\]
ここでさらに、
\[
\sqrt{(a_1+a_2)(a_3+a_4)} \geqq \sqrt{2\sqrt{a_1 a_2} \cdot 2\sqrt{a_3 a_4}} = 2 \sqrt[4]{a_1 a_2 a_3 a_4} \quad \cdots \text{②}
\]
だから、①, ②から
\[
\frac{1}{4}(a_1+a_2+a_3+a_4) \geqq \sqrt[4]{a_1 a_2 a_3 a_4}_{\text{\fbox{同}}}
\]

(ロ) $\frac{a_k}{b_k} > 0$ のため、(1) で $a_k$ に $\frac{a_k}{b_k}$ を代入して、
\[
\sum_{k=1}^4 \frac{a_k}{b_k} \geqq 4 \sqrt[4]{\frac{a_1 a_2 a_3 a_4}{b_1 b_2 b_3 b_4}} = 4 \quad (\because \{a_1, \dots, a_4\} = \{b_1, b_2, b_3, b_4\})
\]
よって示された。$_{\text{\fbox{同}}}$

\end{document}
